\documentclass[11pt]{article}

\usepackage[margin=1in]{geometry}
\usepackage{amsmath,amssymb,amsthm}
\usepackage{enumitem}
\usepackage{xcolor}
\usepackage{hyperref}

\hypersetup{
  colorlinks=true,
  linkcolor=blue!60!black,
  urlcolor=blue!60!black
}

\newcommand{\R}{\mathbb{R}}
\newcommand{\Z}{\mathbb{Z}}
\newcommand{\N}{\mathbb{N}}
\newcommand{\B}{\mathcal{B}}
\newcommand{\F}{\mathcal{F}}
\newcommand{\I}{\mathcal{I}}
\newcommand{\Prob}{\mathbb{P}}
\newcommand{\E}{\mathbb{E}}
\newcommand{\1}{\mathbf{1}}
\newcommand{\Leb}{\lambda}
\newcommand{\Knowledge}[1]{%
  \par\smallskip\noindent
  \textbf{Knowledge used.}
  {\small\color{blue!60!black}\raggedright #1\par}
  \par\smallskip
}

\title{Chapter 6 Exercises}
\author{\textsc{Solutions to Rick Durrett's}\\
\textit{Probability: Theory and Examples}}
\date{\today}

\begin{document}
\maketitle

\noindent
These solutions use the local Chapter 6 LLM/Viki knowledge set in
\path{Probability/LLM_Viki_Dataset/Chapter6_Knowledge}.  The cited
knowledge IDs refer to entries in
\path{chapter6_knowledge_pieces.jsonl}.

\section*{Exercise 6.1.1}

\noindent\textbf{Question.}
Show that the class of invariant events $\I$ is a $\sigma$-field, and
$X\in\I$ if and only if $X$ is invariant, i.e. $X\circ\phi=X$ a.s.

\Knowledge{
\path{durrett_ch6_invariant_sigma_field};
\path{durrett_ch6_ergodicity_definition};
\path{durrett_ch6_measure_preserving_shift_model}.
}

\noindent\textbf{Solution.}
Recall that an event $A\in\F$ is invariant when
$\phi^{-1}A=A$ up to a null set.  First, $\Omega$ and $\varnothing$ are
invariant.  If $A$ is invariant, then
\[
  \phi^{-1}(A^c)=(\phi^{-1}A)^c=A^c
  \quad\text{a.s.},
\]
so $A^c$ is invariant.  If $A_1,A_2,\ldots$ are invariant, then
\[
  \phi^{-1}\left(\bigcup_{j=1}^{\infty}A_j\right)
  =\bigcup_{j=1}^{\infty}\phi^{-1}A_j
  =\bigcup_{j=1}^{\infty}A_j
  \quad\text{a.s.}
\]
Thus $\I$ is a $\sigma$-field.

Here $X\in\I$ means that $X$ is $\I$-measurable.  If $X$ is
$\I$-measurable, then for every rational $r$ the event
$\{X\le r\}$ is invariant.  Hence
\[
  \{X\circ\phi\le r\}
  =\phi^{-1}\{X\le r\}
  =\{X\le r\}
  \quad\text{a.s.}
\]
Taking the intersection over the countable set $r\in\mathbb{Q}$ gives
$X\circ\phi=X$ a.s.

Conversely, if $X\circ\phi=X$ a.s., then for every Borel set
$B\subseteq\R$,
\[
  \phi^{-1}\{X\in B\}
  =\{X\circ\phi\in B\}
  =\{X\in B\}
  \quad\text{a.s.}
\]
Thus every event generated by $X$ is invariant, and $X$ is
$\I$-measurable.

\section*{Exercise 6.1.2}

\noindent\textbf{Question.}
Some authors call $A$ almost invariant if
$\Prob(A\triangle\phi^{-1}A)=0$ and call $C$ invariant in the strict
sense if $C=\phi^{-1}C$.
\begin{enumerate}[label=(\roman*)]
\item Let $A$ be any set, let
$B=\bigcup_{n=0}^{\infty}\phi^{-n}(A)$.  Show
$\phi^{-1}(B)\subset B$.
\item Let $B$ be any set with $\phi^{-1}(B)\subset B$ and let
$C=\bigcap_{n=0}^{\infty}\phi^{-n}(B)$.  Show that
$\phi^{-1}(C)=C$.
\item Show that $A$ is almost invariant if and only if there is a
$C$ invariant in the strict sense with $\Prob(A\triangle C)=0$.
\end{enumerate}

\Knowledge{
\path{durrett_ch6_invariant_sigma_field};
\path{durrett_ch6_measure_preserving_shift_model}.
}

\noindent\textbf{Solution.}
\begin{enumerate}[label=(\roman*)]
\item We have
\[
  \phi^{-1}B
  =\phi^{-1}\left(\bigcup_{n=0}^{\infty}\phi^{-n}A\right)
  =\bigcup_{n=0}^{\infty}\phi^{-(n+1)}A
  =\bigcup_{n=1}^{\infty}\phi^{-n}A
  \subset B.
\]

\item Since $C=\bigcap_{n=0}^{\infty}\phi^{-n}B$,
\[
  \phi^{-1}C
  =\bigcap_{n=1}^{\infty}\phi^{-n}B.
\]
Because $\phi^{-1}B\subset B$, the intersection over $n\ge1$ is already
contained in $B$.  Therefore
\[
  \bigcap_{n=1}^{\infty}\phi^{-n}B
  =B\cap\bigcap_{n=1}^{\infty}\phi^{-n}B
  =C.
\]

\item If there is a strict invariant set $C$ with
$\Prob(A\triangle C)=0$, then
\[
  A\triangle\phi^{-1}A
  \subset (A\triangle C)\cup(C\triangle\phi^{-1}C)
       \cup(\phi^{-1}C\triangle\phi^{-1}A).
\]
The middle term is empty and the other two terms have probability zero,
so $A$ is almost invariant.

Conversely, suppose $A$ is almost invariant.  Then, by induction and
measure preservation,
\[
  \Prob(A\triangle \phi^{-n}A)=0
  \qquad n\ge0.
\]
Let $B=\bigcup_{n=0}^{\infty}\phi^{-n}A$.  Since $B$ differs from $A$
by a countable union of null sets, $\Prob(A\triangle B)=0$.  By part
(i), $\phi^{-1}B\subset B$.  Define
\[
  C=\bigcap_{n=0}^{\infty}\phi^{-n}B.
\]
Part (ii) gives $\phi^{-1}C=C$ strictly.  Also
$\Prob(B\triangle\phi^{-1}B)=0$, and the sets
$B,\phi^{-1}B,\phi^{-2}B,\ldots$ are decreasing, so
$\Prob(B\setminus C)=0$.  Hence $\Prob(A\triangle C)=0$.
\end{enumerate}

\section*{Exercise 6.1.3}

\noindent\textbf{Question.}
A direct proof of ergodicity of rotation of the circle.
\begin{enumerate}[label=(\roman*)]
\item Show that if $\theta$ is irrational, $x_n=n\theta \bmod 1$ is
dense in $[0,1)$.  Hint: all the $x_n$ are distinct, so for any
$N<\infty$, $|x_n-x_m|\le 1/N$ for some $m<n\le N$.
\item Use Exercise A.2.1 to show that if $A$ is a Borel set with
$|A|>0$, then for any $\delta>0$ there is an interval $J=[a,b)$ so that
$|A\cap J|>(1-\delta)|J|$.
\item Combine this with (i) to conclude $\Prob(A)=1$.
\end{enumerate}

\Knowledge{
\path{durrett_ch6_circle_rotation_ergodicity};
\path{durrett_ch6_ergodicity_definition};
\path{durrett_ch6_invariant_sigma_field}.
}

\noindent\textbf{Solution.}
Let $T(x)=x+\theta \bmod 1$.

\begin{enumerate}[label=(\roman*)]
\item The points $0,\theta,2\theta,\ldots,N\theta$ are distinct modulo
1; otherwise $(n-m)\theta$ would be an integer for some $m<n$, which is
impossible when $\theta$ is irrational.  Divide the circle into $N$
intervals of length $1/N$.  Two of the $N+1$ points must fall in the
same interval, so for some $0\le m<n\le N$,
\[
  \|(n-m)\theta\|_{\mathbb{T}}\le \frac1N,
\]
where $\|\cdot\|_{\mathbb{T}}$ is distance to the nearest integer.
Thus the orbit has a nonzero step $q=(n-m)\theta\bmod1$ arbitrarily
close to $0$.  Repeated additions of such a small step move around the
circle with gaps at most $1/N$.  Since $N$ is arbitrary, the orbit
$\{n\theta\bmod1:n\ge0\}$ is dense in $[0,1)$.

\item Exercise A.2.1 gives regular approximation of Borel sets by
finite unions of intervals.  If $|A|>0$, choose a finite union of
intervals $G$ with $|A\triangle G|<\delta |A|/4$.  At least one
component interval $J$ of $G$ must have high relative density of $A$;
otherwise $|G\setminus A|$ would be too large.  Shrinking constants if
necessary, this gives an interval $J=[a,b)$ such that
\[
  |A\cap J|>(1-\delta)|J|.
\]

\item Let $A$ be an invariant Borel set for the irrational rotation and
assume $|A|>0$.  We prove $|A|=1$.  Suppose instead that $|A^c|>0$.
By part (ii), for any small $\delta>0$ there are intervals $J$ and $K$
such that
\[
  |A\cap J|>(1-\delta)|J|,
  \qquad
  |A^c\cap K|>(1-\delta)|K|.
\]
By part (i), the set $\{n\theta\bmod1:n\ge0\}$ is dense, so we may
choose $n$ such that $T^nJ$ overlaps $K$ in an interval $L$ of positive
length.  Choosing $\delta$ small enough compared with $|L|$, the two
sets $T^n(A\cap J)$ and $A^c\cap K$ must have positive intersection
inside $L$.

But invariance of $A$ gives $T^nA=A$ up to null sets.  Hence
$T^n(A\cap J)\subset A$ up to null sets, while $A^c\cap K\subset A^c$.
They cannot have positive intersection, a contradiction.  Therefore
$|A^c|=0$, and $|A|=1$.  The invariant Borel sets have measure $0$ or
$1$, so the irrational rotation is ergodic.
\end{enumerate}

\section*{Exercise 6.1.4}

\noindent\textbf{Question.}
Any stationary sequence $\{X_n,n\ge0\}$ can be embedded in a two-sided
stationary sequence $\{Y_n:n\in\Z\}$.

\Knowledge{
\path{durrett_ch6_stationary_sequence_definition};
\path{durrett_ch6_measure_preserving_shift_model}.
}

\noindent\textbf{Solution.}
We define consistent finite-dimensional distributions on coordinates
indexed by $\Z$.  If $i_1<\cdots<i_k$ are integers, choose $r$ so large
that $i_1+r\ge0$, and set
\[
  \mu_{i_1,\ldots,i_k}
  =\mathcal{L}(X_{i_1+r},\ldots,X_{i_k+r}).
\]
This definition is independent of the chosen $r$, because stationarity
of the one-sided sequence says that shifting all indices forward does
not change the joint law.

The family $\{\mu_{i_1,\ldots,i_k}\}$ is consistent under deletion of
coordinates, again because it comes from finite-dimensional
distributions of the same process after a common shift.  By the
Kolmogorov extension theorem there is a probability measure on the
sequence space $S^{\Z}$ with these finite-dimensional distributions.
Let $Y_n(\omega)=\omega_n$ be the coordinate maps.

For any shift $h\in\Z$ and any finite set of indices,
\[
  \mathcal{L}(Y_{i_1+h},\ldots,Y_{i_k+h})
  =\mu_{i_1+h,\ldots,i_k+h}
  =\mu_{i_1,\ldots,i_k}.
\]
Thus $\{Y_n:n\in\Z\}$ is stationary.  Taking only the coordinates
$n\ge0$ gives the same finite-dimensional distributions as the original
sequence $\{X_n:n\ge0\}$.

\section*{Exercise 6.1.5}

\noindent\textbf{Question.}
If $X_0,X_1,\ldots$ is a stationary sequence and
$g:\R^{\{0,1,\ldots\}}\to\R$ is measurable, then
\[
  Y_k=g(X_k,X_{k+1},\ldots)
\]
is a stationary sequence.  If $X_n$ is ergodic then so is $Y_n$.

\Knowledge{
\path{durrett_ch6_functions_of_stationary_processes};
\path{durrett_ch6_stationary_sequence_definition};
\path{durrett_ch6_ergodicity_definition}.
}

\noindent\textbf{Solution.}
Fix $m\ge0$ and $h\ge0$.  The vector
$(Y_h,\ldots,Y_{h+m})$ is a measurable function of the future sequence
$(X_h,X_{h+1},\ldots)$.  Since $X$ is stationary, the law of
$(X_h,X_{h+1},\ldots)$ is the same as the law of
$(X_0,X_1,\ldots)$.  Therefore
\[
  (Y_h,\ldots,Y_{h+m})
  \stackrel{d}{=}
  (Y_0,\ldots,Y_m),
\]
so $Y$ is stationary.

Now suppose $X$ is ergodic.  Realize $X$ on sequence space with the
shift map $\phi$.  The process $Y$ is a factor of $X$: it is obtained by
applying the same measurable function to each shifted future.  If
$A$ is an invariant event for the $Y$-shift, then the inverse image of
$A$ under the factor map is an invariant event for the $X$-shift.
Since $X$ is ergodic, this inverse image has probability $0$ or $1$.
The event $A$ has the same probability as its inverse image, so $Y$ is
ergodic.

\section*{Exercise 6.1.6}

\noindent\textbf{Question.}
Independent blocks.  Let $X_1,X_2,\ldots$ be a stationary sequence.  Let
$n<\infty$ and let $Y_1,Y_2,\ldots$ be a sequence so that
$(Y_{nk+1},\ldots,Y_{n(k+1)})$, $k\ge0$, are i.i.d. and
\[
  (Y_1,\ldots,Y_n)=(X_1,\ldots,X_n)
  \quad\text{in distribution}.
\]
Finally, let $\nu$ be uniformly distributed on $\{1,2,\ldots,n\}$,
independent of $Y$, and let
\[
  Z_m=Y_{\nu+m},\qquad m\ge1.
\]
Show that $Z$ is stationary and ergodic.

\Knowledge{
\path{durrett_ch6_stationary_sequence_definition};
\path{durrett_ch6_ergodicity_definition};
\path{durrett_ch6_functions_of_stationary_processes};
\path{durrett_ch6_iid_shift_ergodic}.
}

\noindent\textbf{Solution.}
Write the i.i.d. blocks as
\[
  W_k=(Y_{nk+1},\ldots,Y_{n(k+1)}),\qquad k\ge0.
\]
Adjoin an independent phase $\nu$, uniform on $\{1,\ldots,n\}$.  The
sequence $Z$ is what one sees by starting at the random phase $\nu$ in
the concatenation of the i.i.d. blocks.

Stationarity follows from the random phase.  Shifting $Z$ forward by
one coordinate increases the phase by one; if the phase passes $n$, it
wraps back to $1$ and the block sequence is shifted by one block.  The
phase remains uniform, and the shifted block sequence has the same law
because the blocks are i.i.d.  Hence
\[
  (Z_{1+h},\ldots,Z_{m+h})
  \stackrel{d}{=}
  (Z_1,\ldots,Z_m)
\]
for all $m,h\ge1$.

To prove ergodicity, use the same representation.  Let $T$ be the
one-step transformation on ``phase plus i.i.d. blocks.''  The map $T^n$
keeps the phase fixed and shifts the i.i.d. block sequence by one block.
For each fixed phase, the block shift is ergodic by the iid-shift
argument.  If $A$ is invariant under $T$, then it is invariant under
$T^n$, so on each phase section it has probability either $0$ or $1$.
But $T$ carries phase $r$ to phase $r+1$ modulo $n$, so these section
probabilities must be the same for every phase.  Therefore
$\Prob(A)\in\{0,1\}$.  Thus $Z$ is ergodic.

\section*{Exercise 6.1.7}

\noindent\textbf{Question.}
Continued fractions.  Let $\phi(x)=1/x-[1/x]$ for $x\in(0,1)$ and
$A(x)=[1/x]$, where $[1/x]$ is the largest integer $\le 1/x$.
$a_n=A(\phi^n x)$, $n=0,1,2,\ldots$, gives the continued fraction
representation of $x$, i.e.
\[
  x=\frac{1}{a_0+\frac{1}{a_1+\frac{1}{a_2+\cdots}}}.
\]
Show that $\phi$ preserves
\[
  \mu(A)=\frac{1}{\log 2}\int_A\frac{dx}{1+x},
  \qquad A\subset(0,1).
\]

\Knowledge{
\path{durrett_ch6_measure_preserving_shift_model};
\path{durrett_ch6_stationary_sequence_definition};
\path{durrett_ch6_invariant_sigma_field}.
}

\noindent\textbf{Solution.}
It is enough to show $\mu(\phi^{-1}A)=\mu(A)$ for Borel
$A\subset(0,1)$.  For $k\ge1$, the branch on which $[1/x]=k$ is
$x\in(1/(k+1),1/k]$, and
\[
  \phi(x)=\frac1x-k.
\]
Solving $y=\phi(x)$ gives $x=1/(k+y)$, with derivative
$|dx/dy|=(k+y)^{-2}$.  Hence
\[
\begin{aligned}
  \mu(\phi^{-1}A)
  &=\frac{1}{\log2}
    \sum_{k=1}^{\infty}
    \int_A
    \frac{1}{1+\frac{1}{k+y}}\,
    \frac{dy}{(k+y)^2}  \\
  &=\frac{1}{\log2}
    \int_A
    \sum_{k=1}^{\infty}
    \frac{1}{(k+y)(k+y+1)}\,dy.
\end{aligned}
\]
The sum telescopes:
\[
  \sum_{k=1}^{\infty}
  \frac{1}{(k+y)(k+y+1)}
  =
  \sum_{k=1}^{\infty}
  \left(\frac{1}{k+y}-\frac{1}{k+y+1}\right)
  =\frac{1}{1+y}.
\]
Therefore
\[
  \mu(\phi^{-1}A)
  =\frac{1}{\log2}\int_A\frac{dy}{1+y}
  =\mu(A).
\]
Thus the continued-fraction map preserves $\mu$.

\section*{Exercise 6.2.1}

\noindent\textbf{Question.}
Show that if $X\in L^p$ with $p>1$, then the convergence in Theorem
6.2.1 occurs in $L^p$.

\Knowledge{
\path{durrett_ch6_birkhoff_ergodic_theorem};
\path{durrett_ch6_lp_upgrade_birkhoff};
\path{durrett_ch6_invariant_sigma_field}.
}

\noindent\textbf{Solution.}
Let
\[
  A_nX=\frac1n\sum_{m=0}^{n-1}X\circ\phi^m.
\]
Theorem 6.2.1 gives $A_nX\to \E(X\mid\I)$ a.s. and in $L^1$.  We now
upgrade this to $L^p$.

First suppose $X$ is bounded.  Then $A_nX$ is bounded by
$\|X\|_\infty$, and $\E(X\mid\I)$ is also bounded by $\|X\|_\infty$.
Since $A_nX\to\E(X\mid\I)$ a.s., bounded convergence gives
\[
  \E\left|A_nX-\E(X\mid\I)\right|^p\to0.
\]

For general $X\in L^p$, truncate
\[
  X_M=X\mathbf{1}_{\{|X|\le M\}},\qquad R_M=X-X_M.
\]
The bounded case gives
\[
  A_nX_M\to \E(X_M\mid\I)\quad\text{in }L^p.
\]
For the tail, Jensen's inequality gives
\[
  |A_nR_M|^p
  \le \frac1n\sum_{m=0}^{n-1}|R_M\circ\phi^m|^p.
\]
Taking expectations and using measure preservation,
\[
  \|A_nR_M\|_p^p\le \E|R_M|^p.
\]
Also, conditional expectation is a contraction on $L^p$, so
\[
  \|\E(R_M\mid\I)\|_p\le \|R_M\|_p.
\]
Therefore
\[
  \limsup_{n\to\infty}
  \|A_nX-\E(X\mid\I)\|_p
  \le 2\|R_M\|_p.
\]
As $M\to\infty$, $\|R_M\|_p\to0$ by dominated convergence.  Hence
$A_nX\to\E(X\mid\I)$ in $L^p$.

\section*{Exercise 6.2.2}

\noindent\textbf{Question.}
\begin{enumerate}[label=(\roman*)]
\item Show that if $g_n(\omega)\to g(\omega)$ a.s. and
$\E(\sup_k |g_k(\omega)|)<\infty$, then
\[
  \lim_{n\to\infty}\frac1n\sum_{m=0}^{n-1}g_m(\phi^m\omega)
  =\E(g\mid\I)
  \quad\text{a.s.}
\]
\item Show that if we suppose only that $g_n\to g$ in $L^1$, we get
$L^1$ convergence.
\end{enumerate}

\Knowledge{
\path{durrett_ch6_birkhoff_ergodic_theorem};
\path{durrett_ch6_invariant_sigma_field};
\path{durrett_ch6_measure_preserving_shift_model}.
}

\noindent\textbf{Solution.}
\begin{enumerate}[label=(\roman*)]
\item Let $H=\sup_k|g_k|$.  The assumption $\E H<\infty$ implies
$g\in L^1$ and $|g_m-g|\le 2H$.  By Birkhoff's theorem,
\[
  \frac1n\sum_{m=0}^{n-1}g(\phi^m\omega)\to \E(g\mid\I)
  \quad\text{a.s.}
\]
It remains to show that
\[
  \frac1n\sum_{m=0}^{n-1}(g_m-g)(\phi^m\omega)\to0
  \quad\text{a.s.}
\]

For $M\ge1$, define
\[
  q_M(\omega)=\sup_{j\ge M}|g_j(\omega)-g(\omega)|.
\]
Then $q_M\downarrow0$ a.s. and $q_M\le2H$, so
$\E(q_M\mid\I)\downarrow0$ a.s.  For $n>M$,
\[
\begin{aligned}
  \frac1n\sum_{m=0}^{n-1}|g_m-g|(\phi^m\omega)
  &\le
  \frac1n\sum_{m=0}^{M-1}|g_m-g|(\phi^m\omega)
  +\frac1n\sum_{m=M}^{n-1}q_M(\phi^m\omega).
\end{aligned}
\]
The first term tends to $0$ for fixed $M$.  By Birkhoff's theorem, the
limsup of the second term is at most $\E(q_M\mid\I)$ a.s.  Letting
$M\to\infty$ gives the desired a.s. convergence.

\item Again decompose
\[
  \frac1n\sum_{m=0}^{n-1}g_m(\phi^m)
  -\E(g\mid\I)
  =
  \left(\frac1n\sum_{m=0}^{n-1}g(\phi^m)-\E(g\mid\I)\right)
  +\frac1n\sum_{m=0}^{n-1}(g_m-g)(\phi^m).
\]
The first term tends to $0$ in $L^1$ by Theorem 6.2.1.  For the second,
measure preservation gives
\[
  \E\left|
  \frac1n\sum_{m=0}^{n-1}(g_m-g)(\phi^m)
  \right|
  \le
  \frac1n\sum_{m=0}^{n-1}\E|g_m-g|.
\]
Since $\E|g_m-g|\to0$, its Cesaro averages tend to $0$.  Thus the
whole expression converges to $\E(g\mid\I)$ in $L^1$.
\end{enumerate}

\section*{Exercise 6.2.3}

\noindent\textbf{Question.}
Wiener's maximal inequality.  Let $X_j(\omega)=X(\phi^j\omega)$,
$S_k(\omega)=X_0(\omega)+\cdots+X_{k-1}(\omega)$,
$A_k(\omega)=S_k(\omega)/k$, and
$D_k=\max(A_1,\ldots,A_k)$.  Use Lemma 6.2.2 to show that if
$\alpha>0$, then
\[
  \Prob(D_k>\alpha)\le \alpha^{-1}\E|X|.
\]

\Knowledge{
\path{durrett_ch6_maximal_ergodic_lemma};
\path{durrett_ch6_birkhoff_ergodic_theorem}.
}

\noindent\textbf{Solution.}
Apply the maximal ergodic lemma to the integrable random variable
\[
  Y=X-\alpha.
\]
The partial sums of $Y$ along the orbit are
\[
  T_j=S_j-j\alpha,\qquad 1\le j\le k.
\]
Let
\[
  M_k^Y=\max(0,T_1,\ldots,T_k).
\]
Then
\[
  \{M_k^Y>0\}
  =
  \left\{\max_{1\le j\le k}\frac{S_j}{j}>\alpha\right\}
  =
  \{D_k>\alpha\}.
\]
Lemma 6.2.2 gives
\[
  \E(Y;M_k^Y>0)\ge0.
\]
Thus
\[
  0
  \le \E(X-\alpha;D_k>\alpha)
  =\E(X;D_k>\alpha)-\alpha\Prob(D_k>\alpha).
\]
Rearranging,
\[
  \alpha\Prob(D_k>\alpha)
  \le \E(X;D_k>\alpha)
  \le \E(X^+)
  \le \E|X|.
\]
Dividing by $\alpha$ proves
\[
  \Prob(D_k>\alpha)\le \alpha^{-1}\E|X|.
\]

\section*{Exercise 6.3.1}

\noindent\textbf{Question.}
Let
\[
  g_n=\Prob(S_1\ne0,\ldots,S_n\ne0),\qquad n\ge1,
  \qquad g_0=1.
\]
Show that
\[
  \E R_n=\sum_{m=1}^n g_{m-1}.
\]

\Knowledge{
\path{durrett_ch6_range_growth_stationary_increments};
\path{durrett_ch6_birkhoff_ergodic_theorem}.
}

\noindent\textbf{Solution.}
For each point visited by the walk among $S_1,\ldots,S_n$, choose its
last visiting time in $\{1,\ldots,n\}$.  Thus
\[
  R_n=\sum_{m=1}^n
  \mathbf{1}\{S_m\ne S_{m+1},S_m\ne S_{m+2},\ldots,S_m\ne S_n\}.
\]
The event in the $m$th indicator is
\[
  \{S_{m+j}-S_m\ne0 \text{ for }1\le j\le n-m\}.
\]
Since the increments are stationary,
\[
  (S_{m+1}-S_m,\ldots,S_n-S_m)
  \stackrel{d}{=}
  (S_1,\ldots,S_{n-m}).
\]
Therefore
\[
  \Prob(S_m\text{ is not revisited through time }n)=g_{n-m}.
\]
Taking expectations gives
\[
  \E R_n
  =\sum_{m=1}^n g_{n-m}
  =\sum_{r=0}^{n-1}g_r
  =\sum_{m=1}^n g_{m-1}.
\]

\section*{Exercise 6.3.2}

\noindent\textbf{Question.}
Imitate the proof of part (i) in Theorem 6.3.2 to show that if we
assume $\Prob(X_i>1)=0$, $\E X_i>0$, and the sequence $X_i$ is ergodic
in addition to the hypotheses of Theorem 6.3.2, then
\[
  \Prob(A)=\E X_i.
\]

\Knowledge{
\path{durrett_ch6_zero_drift_integer_recurrence};
\path{durrett_ch6_range_growth_stationary_increments};
\path{durrett_ch6_birkhoff_ergodic_theorem}.
}

\noindent\textbf{Solution.}
Let $\mu=\E X_i>0$.  Since the sequence is ergodic, Birkhoff's theorem
gives
\[
  \frac{S_n}{n}\to\mu \qquad\text{a.s.}
\]
Hence
\[
  \frac{\max_{1\le k\le n}S_k}{n}\to\mu.
\]
Indeed, the lower bound follows from $S_n/n\to\mu$; for the upper
bound, fix $K$ so that $S_k/k\le\mu+\varepsilon$ for all $k\ge K$.
Then
\[
  \max_{1\le k\le n}S_k
  \le \max_{1\le k<K}S_k+(\mu+\varepsilon)n,
\]
and let $n\to\infty$, then $\varepsilon\downarrow0$.

Also $\min_{1\le k\le n}S_k/n\to0$.  This is because, eventually,
$S_k\ge(\mu/2)k>0$, so the running minimum is determined by only
finitely many early values.

Since the walk is integer-valued and $X_i\le1$, when the walk reaches a
new positive maximum it cannot skip any integer level.  Thus the range
contains all levels $1,\ldots,\max_{k\le n}S_k$, and
\[
  R_n\ge \max_{1\le k\le n}S_k.
\]
On the other hand every visited site lies between the running minimum
and maximum, so
\[
  R_n\le 1+\max_{1\le k\le n}S_k-\min_{1\le k\le n}S_k.
\]
Combining the previous limits yields
\[
  \frac{R_n}{n}\to\mu \qquad\text{a.s.}
\]
By Theorem 6.3.1,
\[
  \frac{R_n}{n}\to \E(\mathbf{1}_A\mid\I)\qquad\text{a.s.}
\]
Since the increment sequence is ergodic, $\I$ is trivial, so
$\E(\mathbf{1}_A\mid\I)=\Prob(A)$.  Therefore
\[
  \Prob(A)=\mu=\E X_i.
\]

\section*{Exercise 6.3.3}

\noindent\textbf{Question.}
Show that if $\Prob(X_n\in A\text{ at least once})=1$ and
$A\cap B=\varnothing$, then
\[
  \E\left(
    \sum_{1\le m\le T_1}\mathbf{1}_{\{X_m\in B\}}
    \,\middle|\, X_0\in A
  \right)
  =
  \frac{\Prob(X_0\in B)}{\Prob(X_0\in A)}.
\]
When $A=\{x\}$ and $X_n$ is a Markov chain, this is the ``cycle trick''
for defining a stationary measure.

\Knowledge{
\path{durrett_ch6_stationary_return_times_kac};
\path{durrett_ch6_cycle_trick_stationary_measure};
\path{durrett_ch6_stationary_sequence_definition}.
}

\noindent\textbf{Solution.}
Work with a two-sided stationary version of the process.  Let
$T_1=\inf\{m>0:X_m\in A\}$ when $X_0\in A$.  Then
\[
\begin{aligned}
&\E\left(
    \sum_{1\le m\le T_1}\mathbf{1}_{\{X_m\in B\}}
    \,\middle|\, X_0\in A
  \right)  \\
&\qquad =
\frac{1}{\Prob(X_0\in A)}
\sum_{m=1}^{\infty}
\Prob(X_0\in A,\ X_1,\ldots,X_{m-1}\notin A,\ X_m\in B).
\end{aligned}
\]
Because $A\cap B=\varnothing$, the condition $X_m\in B$ automatically
implies $m<T_1$.

By stationarity, the summand equals
\[
  \Prob(X_{-m}\in A,\ X_{-m+1},\ldots,X_{-1}\notin A,\ X_0\in B).
\]
For $m=1,2,\ldots$, these events are disjoint.  Their union is
$\{X_0\in B\}$ up to a null set: on $\{X_0\in B\}$, the assumed
recurrence to $A$ in a stationary two-sided process gives a last visit
to $A$ before time $0$.  Therefore the infinite sum is
$\Prob(X_0\in B)$.  Hence
\[
  \E\left(
    \sum_{1\le m\le T_1}\mathbf{1}_{\{X_m\in B\}}
    \,\middle|\, X_0\in A
  \right)
  =
  \frac{\Prob(X_0\in B)}{\Prob(X_0\in A)}.
\]

\section*{Exercise 6.3.4}

\noindent\textbf{Question.}
Consider the special case in which $X_n\in\{0,1\}$, and let
$\overline{\Prob}=\Prob(\,\cdot\,|X_0=1)$.  Here $A=\{1\}$ and so
\[
  T_1=\inf\{m>0:X_m=1\}.
\]
Show that
\[
  \Prob(T_1=n)
  =
  \frac{\overline{\Prob}(T_1\ge n)}{\overline{\E}T_1}.
\]
When $t_1,t_2,\ldots$ are i.i.d., this reduces to the formula for the
first waiting time in a stationary renewal process.

\Knowledge{
\path{durrett_ch6_stationary_return_times_kac};
\path{durrett_ch6_stationary_renewal_waiting_time};
\path{durrett_ch6_cycle_trick_stationary_measure}.
}

\noindent\textbf{Solution.}
By Kac's return-time theorem applied to $A=\{1\}$,
\[
  \overline{\E}T_1=\frac{1}{\Prob(X_0=1)}.
\]
Thus it remains to identify the numerator.

For $n\ge1$,
\[
  \{T_1=n\}
  =
  \{X_1=0,\ldots,X_{n-1}=0,\ X_n=1\}
\]
under the original stationary law.  By stationarity this event has the
same probability as
\[
  \{X_{-n+1}=0,\ldots,X_{-1}=0,\ X_0=1\}.
\]
Conditioning on $X_0=1$ gives
\[
  \Prob(T_1=n)
  =
  \Prob(X_0=1)\,
  \overline{\Prob}(X_{-n+1}=0,\ldots,X_{-1}=0).
\]
Under $\overline{\Prob}$, the sequence of return gaps to $1$ is
stationary by Theorem 6.3.3.  Hence the previous return gap and the next
return gap have the same distribution, and
\[
  \overline{\Prob}(X_{-n+1}=0,\ldots,X_{-1}=0)
  =
  \overline{\Prob}(T_1\ge n).
\]
Using $\Prob(X_0=1)=1/\overline{\E}T_1$, we obtain
\[
  \Prob(T_1=n)
  =
  \frac{\overline{\Prob}(T_1\ge n)}{\overline{\E}T_1}.
\]

\section*{Exercise 6.5.1}

\noindent\textbf{Question.}
To show that the convergence in part (a) of Theorem 6.4.1 may occur
arbitrarily slowly, let
\[
  X_{m,m+k}=f(k)\ge0,
\]
where $f(k)/k$ is decreasing, and check that $X_{m,m+k}$ is
subadditive.

\Knowledge{
\path{durrett_ch6_subadditive_ergodic_theorem};
\path{durrett_ch6_fekete_expectation_step}.
}

\noindent\textbf{Solution.}
Define $X_{m,n}=f(n-m)$ for $m<n$.  To check subadditivity, let
$0<m<n$.  Since $f(k)/k$ is decreasing,
\[
  \frac{f(m)}{m}\ge \frac{f(n)}{n},
  \qquad
  \frac{f(n-m)}{n-m}\ge \frac{f(n)}{n}.
\]
Multiplying by $m$ and $n-m$ and adding gives
\[
  f(m)+f(n-m)\ge f(n).
\]
Therefore
\[
  X_{0,m}+X_{m,n}=f(m)+f(n-m)\ge f(n)=X_{0,n},
\]
so the array is subadditive.

In this deterministic example,
\[
  \frac{\E X_{0,n}}{n}=\frac{f(n)}{n}.
\]
Thus the convergence in Theorem 6.4.1(a) is exactly the convergence of
the decreasing sequence $f(n)/n$ to its infimum.  By choosing
$f(n)/n$ to decrease to its limit as slowly as desired, the convergence
in part (a) can be made arbitrarily slow.

\section*{Exercise 6.5.2}

\noindent\textbf{Question.}
Consider the longest common subsequence problem, Example 6.4.4, when
$X_1,X_2,\ldots$ and $Y_1,Y_2,\ldots$ are i.i.d. and take the values
$0$ and $1$ with probability $1/2$ each.
\begin{enumerate}[label=(\alph*)]
\item Compute $\E L_1$ and $\E L_2/2$ to get lower bounds on $\gamma$.
\item Show $\gamma<1$ by computing the expected number of $i$ and $j$
sequences of length $K=\alpha n$ with the desired property.
\end{enumerate}

\Knowledge{
\path{durrett_ch6_longest_common_subsequence_limit};
\path{durrett_ch6_subadditive_ergodic_theorem}.
}

\noindent\textbf{Solution.}
\begin{enumerate}[label=(\alph*)]
\item For $n=1$, the longest common subsequence has length $1$ exactly
when $X_1=Y_1$, so
\[
  \E L_1=\Prob(X_1=Y_1)=\frac12.
\]

For $n=2$, $L_2=2$ exactly when the two binary strings are identical,
which has probability $1/4$.  Also $L_2=0$ exactly when one string is
$00$ and the other is $11$, which has probability $2/16=1/8$.
In all remaining cases $L_2=1$.  Hence
\[
  \E L_2
  =2\cdot\frac14+1\cdot\left(1-\frac14-\frac18\right)
  =\frac12+\frac58
  =\frac98.
\]
Therefore
\[
  \gamma\ge \E L_1=\frac12,
  \qquad
  \gamma\ge \frac{\E L_2}{2}=\frac{9}{16}.
\]

\item Let $N_K$ be the number of pairs of increasing index sequences
\[
  1\le i_1<\cdots<i_K\le n,\qquad
  1\le j_1<\cdots<j_K\le n
\]
such that
\[
  X_{i_r}=Y_{j_r},\qquad 1\le r\le K.
\]
For each fixed pair of index sequences, this matching event has
probability $2^{-K}$.  Thus
\[
  \E N_K=\binom{n}{K}^2 2^{-K}.
\]
If $L_n\ge K$, then $N_K\ge1$, so
\[
  \Prob(L_n\ge K)\le \E N_K.
\]

Take $K=\lfloor \alpha n\rfloor$.  Using the entropy estimate
\[
  \binom{n}{\alpha n}
  =\exp\{nH(\alpha)+o(n)\},
  \qquad
  H(\alpha)=-\alpha\log\alpha-(1-\alpha)\log(1-\alpha),
\]
we get
\[
  \E N_K
  =
  \exp\{n(2H(\alpha)-\alpha\log2)+o(n)\}.
\]
For $\alpha<1$ sufficiently close to $1$,
\[
  2H(\alpha)-\alpha\log2<0.
\]
Therefore $\Prob(L_n\ge \alpha n)$ tends to $0$ exponentially fast.
Since $L_n/n\to\gamma$ almost surely by Example 6.4.4, it follows that
$\gamma\le\alpha$ for some $\alpha<1$.  Hence $\gamma<1$.
\end{enumerate}

\section*{Exercise 6.5.3}

\noindent\textbf{Question.}
Given a rate one Poisson process in $[0,\infty)\times[0,\infty)$, let
$(X_1,Y_1)$ be the point that minimizes $x+y$.  Let $(X_2,Y_2)$ be the
point in $[X_1,\infty)\times[Y_1,\infty)$ that minimizes $x+y$, and so
on.  Use this construction to show that in Example 6.5.2
\[
  \gamma\ge (8/\pi)^{1/2}>1.59.
\]

\Knowledge{
\path{durrett_ch6_increasing_subsequence_permutation};
\path{durrett_ch6_poisson_square_time_change};
\path{durrett_ch6_subadditive_ergodic_theorem}.
}

\noindent\textbf{Solution.}
By translation invariance and the memoryless property of the Poisson
process, the increments
\[
  (\Delta X_r,\Delta Y_r)
  =(X_r-X_{r-1},Y_r-Y_{r-1}),\qquad r\ge1,
\]
with $(X_0,Y_0)=(0,0)$, are i.i.d.  Let
$R=\Delta X+\Delta Y$ for one increment.  The event $\{R>t\}$ says that
there is no Poisson point in the triangle
\[
  \{(x,y):x\ge0,\ y\ge0,\ x+y\le t\},
\]
whose area is $t^2/2$.  Hence
\[
  \Prob(R>t)=e^{-t^2/2}.
\]
Therefore
\[
  \E R=\int_0^\infty e^{-t^2/2}\,dt=\sqrt{\frac{\pi}{2}}.
\]
By symmetry, $\E\Delta X=\E\Delta Y=\E R/2=\sqrt{\pi/8}$.

The strong law gives
\[
  \frac{X_k}{k}\to\sqrt{\frac{\pi}{8}},
  \qquad
  \frac{Y_k}{k}\to\sqrt{\frac{\pi}{8}}
  \quad\text{a.s.}
\]
Thus for every $\varepsilon>0$, if
\[
  k\le \frac{t}{\sqrt{\pi/8}+\varepsilon},
\]
then for all large $t$ the first $k$ greedy points lie in the square
$[0,t]\times[0,t]$.  These points form an increasing path, so
\[
  \liminf_{t\to\infty}\frac{Y_{0,t}}{t}
  \ge \frac{1}{\sqrt{\pi/8}+\varepsilon}.
\]
Letting $\varepsilon\downarrow0$ gives
\[
  \gamma\ge \sqrt{\frac{8}{\pi}}>1.59.
\]

\section*{Exercise 6.5.4}

\noindent\textbf{Question.}
Let $\pi_n$ be a random permutation of $\{1,\ldots,n\}$ and let
$J_k^n$ be the number of subsets of $\{1,\ldots,n\}$ of size $k$ so
that the associated $\pi_n(j)$ form an increasing subsequence.  Compute
$\E J_k^n$ and take $k\sim\alpha n^{1/2}$ to conclude that in Example
6.5.2, $\gamma\le e$.

\Knowledge{
\path{durrett_ch6_increasing_subsequence_permutation};
\path{durrett_ch6_poisson_square_time_change}.
}

\noindent\textbf{Solution.}
For each fixed subset $\{j_1<\cdots<j_k\}$, the relative order of
\[
  \pi_n(j_1),\ldots,\pi_n(j_k)
\]
is uniformly distributed over the $k!$ possible orders.  Exactly one of
these orders is increasing, so
\[
  \Prob(\pi_n(j_1)<\cdots<\pi_n(j_k))=\frac1{k!}.
\]
There are $\binom{n}{k}$ subsets, hence
\[
  \E J_k^n=\binom{n}{k}\frac1{k!}.
\]

If $\ell(\pi_n)\ge k$, then $J_k^n\ge1$, so
\[
  \Prob(\ell(\pi_n)\ge k)\le \E J_k^n
  =\binom{n}{k}\frac1{k!}
  \le \frac{n^k}{(k!)^2}.
\]
Let $k\sim \alpha n^{1/2}$.  Stirling's formula gives
\[
  \log\left(\frac{n^k}{(k!)^2}\right)
  =
  k\log n-2k\log k+2k+o(k)
  =
  2\alpha(1-\log\alpha)n^{1/2}+o(n^{1/2}).
\]
If $\alpha>e$, then $1-\log\alpha<0$, and the right-hand side tends to
$-\infty$.  Thus
\[
  \Prob(\ell(\pi_n)\ge \alpha n^{1/2})\to0
  \qquad(\alpha>e).
\]
Since Example 6.5.2 shows that $n^{-1/2}\ell(\pi_n)\to\gamma$, this
implies $\gamma\le \alpha$ for every $\alpha>e$.  Therefore
\[
  \gamma\le e.
\]

\section*{Exercise 6.5.5}

\noindent\textbf{Question.}
Let $\varphi(\theta)=\E e^{-\theta t_i}$ and
\[
  Y_n=(\mu\varphi(\theta))^{-n}
  \sum_{i=1}^{Z_n}e^{-\theta T_n(i)},
\]
where the sum is over individuals in generation $n$ and $T_n(i)$ is the
$i$th person's birth time.  Show that $Y_n$ is a nonnegative martingale
and use this to conclude that if
\[
  \frac{e^{-\theta a}}{\mu\varphi(\theta)}>1,
\]
then
\[
  \Prob(X_{0,n}\le an)\to0.
\]

\Knowledge{
\path{durrett_ch6_age_dependent_branching_speed};
\path{durrett_ch6_branching_speed_large_deviation_formula};
\path{durrett_ch6_subadditive_ergodic_theorem}.
}

\noindent\textbf{Solution.}
Let $\mathcal{F}_n$ be the information up to generation $n$.  Conditional
on $\mathcal{F}_n$, an individual $i$ in generation $n$ has birth time
$T_n(i)$, then lives for a time $t_i$ and produces a random number
$N_i$ of children with mean $\mu$.  These variables are independent of
$\mathcal{F}_n$, and the children are born at time $T_n(i)+t_i$.
Therefore
\[
\begin{aligned}
  \E\left(
    \sum_{j=1}^{Z_{n+1}}e^{-\theta T_{n+1}(j)}
    \,\middle|\,\mathcal{F}_n
  \right)
  &=
  \sum_{i=1}^{Z_n}
  e^{-\theta T_n(i)}
  \E\left(N_i e^{-\theta t_i}\right)  \\
  &=
  \mu\varphi(\theta)
  \sum_{i=1}^{Z_n}e^{-\theta T_n(i)}.
\end{aligned}
\]
Multiplying by $(\mu\varphi(\theta))^{-(n+1)}$ gives
\[
  \E(Y_{n+1}\mid\mathcal{F}_n)=Y_n.
\]
Since $Y_n\ge0$, it is a nonnegative martingale and $\E Y_n=\E Y_0=1$.

Now let
\[
  c=\frac{e^{-\theta a}}{\mu\varphi(\theta)}.
\]
If $X_{0,n}\le an$, then at least one generation-$n$ individual has
birth time at most $an$, so
\[
  \sum_{i=1}^{Z_n}e^{-\theta T_n(i)}\ge e^{-\theta an}.
\]
Hence
\[
  Y_n\ge (\mu\varphi(\theta))^{-n}e^{-\theta an}=c^n.
\]
If $c>1$, Markov's inequality gives
\[
  \Prob(X_{0,n}\le an)
  \le \Prob(Y_n\ge c^n)
  \le c^{-n}\E Y_n
  =c^{-n}\to0.
\]

Equivalently, the condition $e^{-\theta a}/(\mu\varphi(\theta))>1$
means
\[
  \log\mu+\log\varphi(\theta)+\theta a<0,
\]
which is the exponential-moment version of the large-deviation bound
used to identify the speed in the age-dependent branching example.

\end{document}
